\begin{center}
    \section*{\fontsize{20}{20}\selectfont Chapter 6}
\end{center}
\vspace{10mm}

\section{Conclusion}

Our aim was to build a potato leaf disease detector that is \textbf{accurate, small, fast and explainable}, and to make it usable by farmers through a live web application. We used a 5-class field dataset of 2,321 images (Bacteria, Fungi, Healthy, Phytopthora and Virus). We grouped near-duplicate photos with perceptual hashing ($\text{dHash} \le 10$, 2,283 groups) and split the data into 1,859 training, 231 validation and 231 test images. We then tuned a small \textbf{MobileNetV2} model step by step and compared it with three other models on two NVIDIA Tesla T4 GPUs.

The results show that, on a real field dataset, careful training was more useful than a bigger model. We used MobileNetV2 with a small 6,405-parameter head. We restored the best Stage~1 weights, kept Batch Normalization frozen in the top 54 layers, used AdamW and added leaf-like augmentation. With this setup our \textbf{8.64~MB (2.26M-parameter)} model reached \textbf{93.51\% test accuracy}, \textbf{0.9272 macro-F1} and \textbf{0.4466 test loss}. With 4-view Test-Time Augmentation, which adds no parameters, it reached \textbf{93.94\% accuracy} and \textbf{0.9317 macro-F1}. On the same test images it was both smaller and more accurate than every other model. It beat the MultiNet-V3 ensemble, which is $13.7\times$ larger (118.55~MB, 90.91\% accuracy, 0.8975 macro-F1), as well as InceptionV3 (85.18~MB, 84.42\%) and NASNetMobile (17.33~MB, 72.73\%). The exported model gave exactly the same output as the trained one ($0.00 \times 10^{0}$ difference).

Grad-CAM on \texttt{Conv\_1} showed that the model looks at the lesions on the leaf. We put the whole pipeline into a live FastAPI and React web application in English and Bengali (\url{https://potato.mamunlab.com/}). It has phone camera capture, batch analysis of 12 images, PDF/CSV reports and Google Gemini chat with safety rules. Together, these parts form a complete farm AI tool that others can rebuild and that is made for field use.

\subsection{Limitations}
The system still has some clear limits.

\textit{Small test set.} The leak-aware split stops duplicates from raising the score, but the test set is only 231 images. It has just 19 \textit{Healthy} and 31 \textit{Phytophthora} leaves. Recall is 100\% for \textit{Healthy} and 94\% for \textit{Phytophthora}, but \textit{Healthy} precision is only 0.79, because a few leaves with mild fungal disease were called healthy.

\textit{One crop and five classes.} The model knows only five potato leaf classes. It cannot yet detect insect damage, nematodes, nutrient problems or diseases of other local crops such as rice, tomato or wheat.

\textit{Needs the internet.} The 8.64~MB CNN is small enough to run on a device, but the Gemini chat and the web app need an internet connection. In remote farm areas the connection is often poor.

\textit{No real field trial.} We tested on a public field dataset. We did not run a trial with farmers over several seasons, with local plant doctors checking the results.

\subsection{Future Works and Direction}
To turn this capstone into a tool that many rural farmers can use, we plan five next steps.

\textit{Offline use on the phone (TFLite / INT8).} INT8 quantization after training could shrink the 8.64~MB model to about 2.5~MB. It could then run fully offline on a phone with TensorFlow Lite (TFLite) or ONNX Runtime Web. The disease cards could be saved on the phone, and the Gemini chat could sync when the internet comes back.

\textit{More data and rejection of unknown images.} We want to collect more potato photos from several seasons, especially mild fungal leaves and healthy leaves. We also want to add a check based on confidence or entropy that rejects photos that are not potato leaves. This should raise \textit{Healthy} precision.

\textit{More crops.} We plan to use the same training recipe (frozen BN, AdamW) and the same leak-aware benchmark for other important local crops, such as rice, wheat, tomato and eggplant.

\textit{Advice based on local manuals.} We want to add Retrieval-Augmented Generation (RAG) to the Gemini module, using checked guidelines from the local agriculture department. The Bengali and English advice would then fit the region better and would be easier to check.

\textit{Field tests with agronomists.} We hope to work with agricultural extension officers to test the Grad-CAM overlays, the TTA predictions and the advice in both languages on real small potato farms during the blight season.
